\chapter{Preliminaries}\label{chap:preliminaries}

The opening paragraph of a chapter should briefly outline the chapter structure 
by naming each section and stating what it covers.

In the preliminaries you can name the formalism studied (e.g., term rewriting 
systems, integer transition systems, probabilistic programs) and cite a 
standard reference for further reading, e.g., \cite{baader_nipkow_1999}. 
Do not add a citation to every definition.
%-----------------------------------------------------------------------
\section{Basic Definitions}\label{sec:basic-definitions}
%-----------------------------------------------------------------------

Introduce the formal objects of the theory bottom-up: start with the most
elementary notion (e.g., the syntactic objects) and build towards the
property of interest (e.g., termination, complexity).

Each definition in this section should be accompanied by:
\begin{enumerate}
    \item A sentence stating what the definition formalizes and why it is needed.
    \item The formal statement in a Definition box.
    \item A brief explanation of each component of the definition.
    \item At least one example.
\end{enumerate}

Cover the following in order:
\begin{enumerate}
    \item \textbf{Syntactic Objects.} Define the elementary building blocks
          of the formalism (e.g., the signature and the set of terms, the
          program syntax, the set of states and transitions).
    \item \textbf{Evaluation Relation.} Define how one object
          evaluates to another. If several strategies exist, introduce them
          here and state which one is used in the rest of the thesis.
    \item \textbf{Derived Notions.} Define derived concepts such as normal
          forms, reachability, derivation lengths, etc.
    \item \textbf{The Property of Interest.} Give the formal definition of
          the property the thesis analyzes (e.g., termination, innermost
          complexity, almost-sure termination).
\end{enumerate}

% --- EXAMPLE: style illustration from a thesis on term rewriting ---
% Replace the definitions below with those relevant to your own topic.

In the begging, we start by defining [Basic Notion] that is used throughout the 
remainder of this thesis.

\begin{Definition}{[Name of First Basic Notion]}{first-notion}
    A [Thing] is [Basic Notion] if it satisfies [Property 1] and [Property 2].
    % Define the most elementary object of the formalism.
    % Example (term rewriting): the signature and the set of terms.
\end{Definition}

The reason for requiring [Property 1] is [...]. 
[Property 2] is generally assumed in the literature.
% Explain the definition in one or two sentences, then give an example.

\begin{example}\label{ex:first-notion}
    % Illustrate the definition on a small, concrete input.
    % If you use the same running example throughout the chapter, introduce
    % it here and mark it as the running example.
\end{example}

\begin{example}\label{ex:first-notion-second-example}
    % You can add several examples for one definition, if the definition 
    % is complex and hard to understand.
\end{example}


%-----------------------------------------------------------------------
\section{Known Results from the Literature}\label{sec:known-results}
%-----------------------------------------------------------------------

Present the lemmas and theorems from the literature that are used in
\cref{chap:main-chapter}. 
You do not have to prove these results unless the proof is insightful for the 
rest of the thesis. 
Always cite the original source.
%
Each result in this section should be presented as follows:
\begin{enumerate}
    \item A sentence explaining why the result is important and how it will
     be used in the thesis.
     \item The formal statement in a Lemma or Theorem box, with a citation in
     the title.
     \item A proof indication: either a short proof, ``Proof: see 
     \cite{replace-with-citation}'', or a forward reference to the appendix.
     \item (Optional) A brief remark on the scope or limitations of the result
     if that is relevant for understanding its use in \cref{chap:main-chapter}.
\end{enumerate}

This section collects known results from the literature that are used in the
rest of the thesis.
Organize them as follows:
\begin{enumerate}
    \item \textbf{Auxiliary lemmas.} State small technical results about the
          basic objects from \Cref{sec:basic-definitions} (e.g., properties
          of substitution, closure conditions, or monotonicity).
          These are typically needed in the proofs of \Cref{chap:main-chapter}.

    \item \textbf{Characterization theorem(s).} State the main theoretical
          results from the literature about the property of interest.
          These serve as a reference point for understanding what is already
          known.

    \item \textbf{Existing automated technique(s).} If the thesis extends an
          existing automated method (e.g., dependency pairs, polynomial
          interpretations, ranking functions), state its correctness theorem
          here.
\end{enumerate}

% --- Example entries (replace with results relevant to your topic) ---

\begin{Lemma}{[Name of Auxiliary Lemma], \cite{giesl2006mechanizing}}{aux-lemma}
    A Lemma is a small result or auxiliary result, which is only interesting because 
    we can use it to prove the main result.
\end{Lemma}

\begin{proof}
    The complete proof can be found in \cite{giesl2006mechanizing}.
\end{proof}

\begin{Theorem}{[Name of Key Theorem], \cite{giesl2006mechanizing}}{key-theorem}
    A Theorem is one of the main results. Everything else is defined and proven such
    that this results can be proven.
\end{Theorem}

\begin{proof}
    The complete proof can be found in \cite{giesl2006mechanizing}.
\end{proof}

\begin{example}
    Apply the theorem above to a concrete example.
    Demonstrate both a case where the result applies successfully and,
    if possible, a case where it fails due to the restriction identified above.
    The failing case motivates the extension in \cref{chap:main-chapter}.
\end{example}

\begin{Corollary}{[Name of a Direct Consequence], \cite{giesl2006mechanizing}}{key-corollary}
    A Corollary is a (easy) consequence of some other (main) result. 
    Still, it has to be proven, even if the proof is rather simple!
\end{Corollary}

\begin{proof}
    The complete proof can be found in \cite{giesl2006mechanizing}.
\end{proof}

\begin{example}

\end{example}

% Add further Lemma and Theorem boxes as needed.

You can use \textcolor{ForestGreen}{colors} to highlight differences in boxes 
or even difficult parts in proofs. 
However, do not make the colors too distracting and everything has to be
readable in black and white as well.